  \documentclass{standalone}
  \usepackage{amsmath}
  \usepackage{amsfonts}
  \usepackage{tikz}
  \usepackage{tikz-3dplot}
  \usepackage{pgfplots}
  \usepackage{pst-plot}
  \pgfplotsset{compat=1.18} % version number
  \usepackage{physics}
  \usepackage[outline]{contour} % glow around text
  \begin{document}
  \begin{tikzpicture}

\pgfplotsset{
    soldot/.style={color=black,only marks,mark=*},
    opendot/.style={color=black,fill=white,only marks,mark=*},
    compat=1.12
}
\begin{axis}[
    axis lines=middle,
    grid=both,
    xlabel=\(x\),
    ylabel=\(y\),
    xmin=-1.5, xmax=4.5,
    ymin=-1, ymax=5,
    xtick=\empty,
    ytick=\empty,
    samples=200,
    width=10cm,
    height=8cm,
    domain=-0.5:4.5,
    ticklabel style={fill=white},
]
    % Parabola: nice clean shape
    \addplot[blue, thick, smooth] {0.5*(x-2)^2 + 1};

    % Points P and Q
    \addplot[soldot] coordinates{(1.5,1.125)} node[below left] {P};
    \addplot[soldot] coordinates{(3,2.5)} node[above right] {Q};

    % Dotted lines to axes from P and Q
    \draw[dashed] (axis cs:1.5,0) -- (axis cs:1.5,1.125);
    \draw[dashed] (axis cs:0,1.125) -- (axis cs:1.5,1.125);
    \draw[dashed] (axis cs:3,0) -- (axis cs:3,2.5);
    \draw[dashed] (axis cs:0,2.5) -- (axis cs:3,2.5);

    % Secant line PQ
    \addplot[domain=1:3.5, dashed] {0.9167*(x-1.5)+1.125};

    % Labels
    \node at (axis cs:1.5,-0.4) {\(x\)};
    \node at (axis cs:3,-0.4) {\(x+\Delta x\)};
    \node at (axis cs:-1.0,1.125) {\(f(x)\)};
    \node at (axis cs:-1.0,2.5) {\(f(x+\Delta x)\)};
    \node at (axis cs:2.25,-0.8) {\(\Delta x\)};
\end{axis}
  \end{tikzpicture}
  \end{document}
